% GENERATED FILE. Edit canonical lessons, then run npm run generate:books.
% canonical-source-hash: fnv1a64:044fecabc1fffc26
% canonical-lessons: TA-C201-nerri, TA-C201-tati, TA-C201-micai, TA-C201-ullankai, TA-C201-vayirruvali

\chapter{Forehead, Beard, Moustache, Palm of the hand, Stomach ache}
\label{ch:ta-forehead-beard-moustache-palm-of-the-hand-stomach-ache}
\hlchaptermodality{201}

\begin{chapteropening}
\textbf{By the end of this chapter:} \emph{I can say the forehead, a beard, a moustache, the palm of the hand and a stomach ache.}

Complete the last lesson of chapter 201: I can say the forehead, a beard, a moustache, the palm of the hand and a stomach ache.
\end{chapteropening}

\section[\texorpdfstring{\ta{நெற்றி} (neṟṟi)}{நெற்றி (neṟṟi)}]{\ta{நெற்றி} (neṟṟi) — the forehead}

\label{lesson:TA-C201-nerri}



\begin{quote}
Before the new one: say the Tamil for to collect, then the Tamil for to suck.
\end{quote}

\subsection*{You'll want to know: \ta{நெற்றி}}
\textbf{\ta{நெற்றி}} — \emph{neṟṟi} — \textquotedblleft{}the forehead\textquotedblright{}.

\begin{grammarlens}[title={What you've built}]
One more word for this chapter.
\end{grammarlens}

\subsection*{Your turn}
\begin{itemize}
\raggedright
  \item \emph{Say it:} \emph{neṟṟi}
  \item \emph{Say it:} \emph{neṟṟi}, once more
  \item \emph{Say it:} say \emph{neṟṟi}
  \item \emph{Recall:} say the Tamil for to collect, then the Tamil for to suck, then say \emph{neṟṟi} again
\end{itemize}

\begin{tcolorbox}[breakable,colback=teal!4,colframe=teal!35!black,title={Before you move on}]
What does \ta{நெற்றி} mean? (\textquotedblleft{}The forehead\textquotedblright{}.) Say it once more.
\end{tcolorbox}
\section[\texorpdfstring{\ta{தாடி} (tāṭi)}{தாடி (tāṭi)}]{\ta{தாடி} (tāṭi) — a beard}

\label{lesson:TA-C201-tati}



\begin{quote}
Before the new one: say the Tamil for to suck, then the Tamil for the forehead.
\end{quote}

\subsection*{You'll want to know: \ta{தாடி}}
\textbf{\ta{தாடி}} — \emph{tāṭi} — \textquotedblleft{}a beard\textquotedblright{}.

\begin{grammarlens}[title={What you've built}]
One more word for this chapter.
\end{grammarlens}

\subsection*{Your turn}
\begin{itemize}
\raggedright
  \item \emph{Say it:} \emph{tāṭi}
  \item \emph{Say it:} \emph{tāṭi}, once more
  \item \emph{Say it:} say \emph{tāṭi}
  \item \emph{Recall:} say the Tamil for to suck, then the Tamil for the forehead, then say \emph{tāṭi} again
\end{itemize}

\begin{tcolorbox}[breakable,colback=teal!4,colframe=teal!35!black,title={Before you move on}]
What does \ta{தாடி} mean? (\textquotedblleft{}A beard\textquotedblright{}.) Say it once more.
\end{tcolorbox}
\section[\texorpdfstring{\ta{மீசை} (mīcai)}{மீசை (mīcai)}]{\ta{மீசை} (mīcai) — a moustache}

\label{lesson:TA-C201-micai}



\begin{quote}
Before the new one: say the Tamil for the forehead, then the Tamil for a beard.
\end{quote}

\subsection*{You'll want to know: \ta{மீசை}}
\textbf{\ta{மீசை}} — \emph{mīcai} — \textquotedblleft{}a moustache\textquotedblright{}.

\begin{grammarlens}[title={What you've built}]
One more word for this chapter.
\end{grammarlens}

\subsection*{Your turn}
\begin{itemize}
\raggedright
  \item \emph{Say it:} \emph{mīcai}
  \item \emph{Say it:} \emph{mīcai}, once more
  \item \emph{Say it:} say \emph{mīcai}
  \item \emph{Recall:} say the Tamil for the forehead, then the Tamil for a beard, then say \emph{mīcai} again
\end{itemize}

\begin{tcolorbox}[breakable,colback=teal!4,colframe=teal!35!black,title={Before you move on}]
What does \ta{மீசை} mean? (\textquotedblleft{}A moustache\textquotedblright{}.) Say it once more.
\end{tcolorbox}
\section[\texorpdfstring{\ta{உள்ளங்கை} (uḷḷaṅkai)}{உள்ளங்கை (uḷḷaṅkai)}]{\ta{உள்ளங்கை} (uḷḷaṅkai) — the palm of the hand}

\label{lesson:TA-C201-ullankai}



\begin{quote}
Before the new one: say the Tamil for a beard, then the Tamil for a moustache.
\end{quote}

\subsection*{You'll want to know: \ta{உள்ளங்கை}}
\textbf{\ta{உள்ளங்கை}} — \emph{uḷḷaṅkai} — \textquotedblleft{}the palm of the hand\textquotedblright{}.

\begin{grammarlens}[title={What you've built}]
One more word for this chapter.
\end{grammarlens}

\subsection*{Your turn}
\begin{itemize}
\raggedright
  \item \emph{Say it:} \emph{uḷḷaṅkai}
  \item \emph{Say it:} \emph{uḷḷaṅkai}, once more
  \item \emph{Say it:} say \emph{uḷḷaṅkai}
  \item \emph{Recall:} say the Tamil for a beard, then the Tamil for a moustache, then say \emph{uḷḷaṅkai} again
\end{itemize}

\begin{tcolorbox}[breakable,colback=teal!4,colframe=teal!35!black,title={Before you move on}]
What does \ta{உள்ளங்கை} mean? (\textquotedblleft{}The palm of the hand\textquotedblright{}.) Say it once more.
\end{tcolorbox}
\section[\texorpdfstring{\ta{வயிற்றுவலி} (vayiṟṟuvali)}{வயிற்றுவலி (vayiṟṟuvali)}]{\ta{வயிற்றுவலி} (vayiṟṟuvali) — a stomach ache}

\label{lesson:TA-C201-vayirruvali}



\begin{quote}
Before the new one: say the Tamil for a moustache, then the Tamil for the palm of the hand.
\end{quote}

\subsection*{You'll want to know: \ta{வயிற்றுவலி}}
\textbf{\ta{வயிற்றுவலி}} — \emph{vayiṟṟuvali} — \textquotedblleft{}a stomach ache\textquotedblright{}.

\begin{grammarlens}[title={What you've built}]
One more word for this chapter.
\end{grammarlens}

\subsection*{Your turn}
\begin{itemize}
\raggedright
  \item \emph{Say it:} \emph{vayiṟṟuvali}
  \item \emph{Say it:} \emph{vayiṟṟuvali}, once more
  \item \emph{Say it:} say \emph{vayiṟṟuvali}
  \item \emph{Recall:} say the Tamil for a moustache, then the Tamil for the palm of the hand, then say \emph{vayiṟṟuvali} again
\end{itemize}

\begin{tcolorbox}[breakable,colback=teal!4,colframe=teal!35!black,title={Before you move on}]
What does \ta{வயிற்றுவலி} mean? (\textquotedblleft{}A stomach ache\textquotedblright{}.) Say it once more.
\end{tcolorbox}
